\chapter{Phases of the Moon}\label{ch:g7:moon-phases}

Years ago you kept a \omterm{def:g2:moon-in-the-sky:moon}{Moon} diary and watched the silver shape march
from \omterm{ex:g2:moon-in-the-sky:shapes}{crescent} to full and back, promised the ``why'' for later.
Later is now. You own everything the explanation needs: a \omterm{def:g2:moon-in-the-sky:moon}{Moon}
that \omterm{def:g6:sun-earth-moon:axisorbit}{orbits} monthly, sunlight in straight lines, and a sphere
that is always --- always --- half lit.

\section{The one fact underneath}

\begin{proposition}[Half-lit, always]\label{prop:g7:moon-phases:halflit}
The Sun \omterm{def:g1:light-and-shadows:source}{lights} one half of the \omterm{def:g2:moon-in-the-sky:moon}{Moon} at every moment --- the half
facing it --- exactly as it \omterm{def:g1:light-and-shadows:source}{lights} half the Earth. The \omterm{def:g2:moon-in-the-sky:moon}{Moon} has
no changing shape, no growing and shrinking: it is a stone ball
wearing a permanent day side and a permanent night side, the
boundary sliding as it \omterm{def:g6:sun-earth-moon:axisorbit}{orbits}. Everything we call a ``\omterm{def:g7:moon-phases:phases}{phase}'' is
merely the \emph{view}: how much of the day side happens to face
the Earth.
\end{proposition}

\begin{definition}[Phases]\label{def:g7:moon-phases:phases}
The \emph{phases}\index{phase} of the \omterm{def:g2:moon-in-the-sky:moon}{Moon} are its apparent
shapes through the month: \emph{new moon} (day side turned fully
away --- we see nothing), \emph{\omterm{ex:g2:moon-in-the-sky:shapes}{crescent}}, \emph{first quarter}
(half the disc lit), \emph{gibbous} (more than half),
\emph{\omterm{ex:g2:moon-in-the-sky:shapes}{full moon}} (day side turned fully toward us), then gibbous,
\emph{last quarter} and \omterm{ex:g2:moon-in-the-sky:shapes}{crescent} again, closing the cycle in
about $29.5$ days. From new to full the lit part grows night by
night: the \omterm{def:g2:moon-in-the-sky:moon}{Moon} is \emph{waxing}\index{waxing}; from full back
to new it shrinks: \emph{waning}\index{waning}.
\end{definition}

\begin{omfigure}
\begin{tikzpicture}[scale=0.8]
  % dashed orbit
  \draw[dashed, omExo!70] (0,0) ellipse (4.2 and 3.1);
  % eight stations: dark half toward the left, sunlit half (white)
  % toward the sunlight arriving from the right
  \foreach \a in {0,45,...,315} {
    \begin{scope}[shift={({4.2*cos(\a)},{3.1*sin(\a)})}]
      \fill[omExo] (0,0) circle (0.46);
      \begin{scope}
        \clip (0,0) circle (0.46);
        \fill[white] (0,-0.5) rectangle (0.5,0.5);
      \end{scope}
      \draw[thick] (0,0) circle (0.46);
      \draw[thick] (0,0.46) -- (0,-0.46);
    \end{scope}
  }
  % the Earth, its own day side facing the sunlight
  \fill[omDef!35] (0,0) circle (0.55);
  \begin{scope}
    \clip (0,0) circle (0.55);
    \fill[omDef!12] (0,-0.6) rectangle (0.6,0.6);
  \end{scope}
  \draw[thick] (0,0) circle (0.55);
  \node[font=\footnotesize] at (0,0) {Earth};
  \foreach \y in {-2.4,-1.2,0,1.2,2.4} {
    \draw[-Stealth, omThm, thick] (7.7,\y) -- (6.4,\y);
  }
  \node[right, font=\small, omThm] at (7.8,0.55) {sunlight};
  \node[left, font=\footnotesize] at (-4.85,0) {full};
  \node[right, font=\footnotesize] at (4.85,0) {new};
  \node[above, font=\footnotesize] at (0,3.75) {last quarter};
  \node[below, font=\footnotesize] at (0,-3.75) {first quarter};
\end{tikzpicture}

{\small The whole secret in one picture: eight stations of the
orbiting \omterm{def:g2:moon-in-the-sky:moon}{Moon}, each half-lit on its sunward side. What changes is
only which slice of the day side faces the Earth.}
\end{omfigure}

\begin{method}[Be the Earth, hold the Moon]\label{met:g7:moon-phases:walk}
The explanation, performed in one minute:
\begin{enumerate}
  \item in a dark room, a single bare lamp is the Sun; you are
        the Earth; a white ball on your outstretched hand is the
        \omterm{def:g2:moon-in-the-sky:moon}{Moon};
  \item face the lamp, ball toward it: the ball's lit half faces
        away from you --- you see its dark side: new moon;
  \item turn slowly on the spot (the ball orbiting you): a lit
        sliver appears at the ball's edge --- \omterm{def:g7:moon-phases:phases}{waxing} \omterm{ex:g2:moon-in-the-sky:shapes}{crescent};
        after a quarter turn, half the ball shows lit: first
        quarter;
  \item keep turning: gibbous --- then, lamp at your back, ball
        fully lit before you: \omterm{ex:g2:moon-in-the-sky:shapes}{full moon}; complete the turn
        through \omterm{def:g7:moon-phases:phases}{waning} half and \omterm{ex:g2:moon-in-the-sky:shapes}{crescent} to new.
\end{enumerate}
One turn of yourself replays one month of sky.
\end{method}

\section{The misconception, retired}

\begin{proposition}[Phases are not the Earth's shadow]\label{prop:g7:moon-phases:notshadow}
The most popular wrong answer in astronomy: ``the \omterm{def:g7:moon-phases:phases}{phases} are the
Earth's \omterm{def:g1:light-and-shadows:shadow}{shadow} on the \omterm{def:g2:moon-in-the-sky:moon}{Moon}.'' Geometry refutes it three ways.
The Earth's \omterm{def:g1:light-and-shadows:shadow}{shadow} points \emph{away} from the Sun --- it can
only touch a \omterm{def:g2:moon-in-the-sky:moon}{Moon} standing opposite the Sun, which happens
solely at \omterm{ex:g2:moon-in-the-sky:shapes}{full moon} (and then only at rare eclipses). A \omterm{ex:g2:moon-in-the-sky:shapes}{crescent}
\omterm{def:g2:moon-in-the-sky:moon}{Moon} stands \emph{sunward} of us, nowhere near our \omterm{def:g1:light-and-shadows:shadow}{shadow}. And
the \omterm{def:g1:light-and-shadows:shadow}{shadow}'s edge on the \omterm{def:g2:moon-in-the-sky:moon}{Moon} is always part of a large circle
--- it could never carve the gibbous shape's double-curved
boundary. The \omterm{def:g7:moon-phases:phases}{phases} need no \omterm{def:g1:light-and-shadows:shadow}{shadow}: a half-lit ball, seen from
different sides, does everything.
\end{proposition}

\begin{example}[Eclipse versus phase]\label{ex:g7:moon-phases:eclipsevs}
Keep the two Moon-darkenings in separate drawers. \emph{\omterm{def:g7:moon-phases:phases}{Phase}}:
the slow monthly cycle, the \omterm{def:g2:moon-in-the-sky:moon}{Moon}'s own night side turning toward
us --- no \omterm{def:g1:light-and-shadows:shadow}{shadow} of ours involved. \emph{\omterm{def:g7:shadows-eclipses:lunar}{Lunar eclipse}}: a rare
hour when the full \omterm{def:g2:moon-in-the-sky:moon}{Moon} crosses our \omterm{prop:g7:shadows-eclipses:umbra}{umbra} --- the Earth's \omterm{def:g1:light-and-shadows:shadow}{shadow}
genuinely falling on it, coppery and swift. The diary's rhythm is
a view; the eclipse is a visit.
\end{example}

\section{Reading the sky like a clock}

\begin{example}[Where the phases keep their hours]\label{ex:g7:moon-phases:hours}
The diagram fixes each \omterm{def:g7:moon-phases:phases}{phase} to a position --- and each position
to a schedule. The full \omterm{def:g2:moon-in-the-sky:moon}{Moon} stands opposite the Sun: it must
rise at \omterm{def:g1:day-and-night:daynight}{sunset}, reign all night, and set at \omterm{def:g1:day-and-night:daynight}{sunrise}. The new
\omterm{def:g2:moon-in-the-sky:moon}{Moon} travels \emph{with} the Sun: up all day, lost in glare ---
which is why nobody ``sees'' a new moon. First quarter trails
the Sun by a quarter turn: it rises near midday and sets near
midnight --- the after-dinner \omterm{def:g2:moon-in-the-sky:moon}{Moon}. The old \omterm{def:g2:moon-in-the-sky:moon}{daytime-Moon} surprise
of your childhood is now a timetable: any \omterm{def:g2:moon-in-the-sky:moon}{Moon} not new nor full
shares part of its shift with the Sun.
\end{example}

\begin{example}[Which way does the crescent point?]\label{ex:g7:moon-phases:points}
The lit side of any \omterm{def:g7:moon-phases:phases}{phase} faces the Sun --- the \omterm{ex:g2:moon-in-the-sky:shapes}{crescent} is a
bow aimed at its archer. Evening \omterm{ex:g2:moon-in-the-sky:shapes}{crescent} in the west: horns
pointing up-and-away from the \omterm{def:g1:day-and-night:daynight}{sunset} glow beneath the horizon.
The sky never draws it wrong; painters often do. A \omterm{ex:g2:moon-in-the-sky:shapes}{crescent} with
its horns toward the Sun is astronomy's favorite spot-the-error.
\end{example}

\begin{remark}[The face that never turns away]\label{rem:g7:moon-phases:sameface}
One more diary observation, finally honored: the \omterm{def:g2:moon-in-the-sky:moon}{Moon}'s dark
patches form the same face every night, \omterm{def:g7:moon-phases:phases}{phase} after \omterm{def:g7:moon-phases:phases}{phase}, year
after year. The \omterm{def:g2:moon-in-the-sky:moon}{Moon} turns exactly once on its \omterm{def:g6:sun-earth-moon:axisorbit}{axis} per \omterm{def:g6:sun-earth-moon:axisorbit}{orbit}
--- so the same half faces the Earth forever. No one saw the far
side until a space probe flew around with a camera within
living memory of your grandparents. Why the turning and the
orbiting agree so perfectly is a deep story of tides, kept for
the university years of this course.
\end{remark}

\section{Exercises}

\begin{exercise}[$\star$]\label{exo:g7:moon-phases:1}
What fraction of the \omterm{def:g2:moon-in-the-sky:moon}{Moon} is sunlit at any moment? What, then,
is a ``\omterm{def:g7:moon-phases:phases}{phase}''?
\end{exercise}

\begin{exercise}[$\star$]\label{exo:g7:moon-phases:2}
Recite the \omterm{def:g7:moon-phases:phases}{phases} in order from new moon, and give the cycle's
length. Which half of the cycle is ``\omterm{def:g7:moon-phases:phases}{waxing}''?
\end{exercise}

\begin{exercise}[$\star$]\label{exo:g7:moon-phases:3}
In the orbit-walk method, where are lamp, ball and you at new
moon? At \omterm{ex:g2:moon-in-the-sky:shapes}{full moon}?
\end{exercise}

\begin{exercise}[$\star$]\label{exo:g7:moon-phases:4}
Give two of the three geometric refutations of ``the \omterm{def:g7:moon-phases:phases}{phases} are
the Earth's \omterm{def:g1:light-and-shadows:shadow}{shadow}''.
\end{exercise}

\begin{exercise}[$\star$]\label{exo:g7:moon-phases:5}
What is the difference between the new moon's darkness and a
\omterm{def:g7:shadows-eclipses:lunar}{lunar eclipse}'s darkness?
\end{exercise}

\begin{exercise}[$\star$]\label{exo:g7:moon-phases:6}
Why does the full \omterm{def:g2:moon-in-the-sky:moon}{Moon} always rise near \omterm{def:g1:day-and-night:daynight}{sunset}? Argue from its
position in the diagram.
\end{exercise}

\begin{exercise}[$\star$]\label{exo:g7:moon-phases:7}
Why does nobody ever ``see the new moon rise at midnight''?
Where is the new moon at midnight?
\end{exercise}

\begin{exercise}[$\star$]\label{exo:g7:moon-phases:8}
An evening \omterm{ex:g2:moon-in-the-sky:shapes}{crescent} hangs in the west. Which way do its horns
point relative to the vanished Sun --- and why can a \omterm{ex:g2:moon-in-the-sky:shapes}{crescent}
never aim its horns at the Sun?
\end{exercise}

\begin{exercise}[$\star\star$]\label{exo:g7:moon-phases:9}
Tonight the \omterm{def:g2:moon-in-the-sky:moon}{Moon} is \omterm{def:g7:moon-phases:phases}{waxing} gibbous. Where in its \omterm{def:g6:sun-earth-moon:axisorbit}{orbit} does it
stand? Roughly when did it rise, and what \omterm{def:g7:moon-phases:phases}{phase} --- and what
schedule --- comes five nights later?
\end{exercise}

\begin{exercise}[$\star\star$]\label{exo:g7:moon-phases:10}
A \omterm{def:g7:shadows-eclipses:solar}{solar eclipse} can only happen at one \omterm{def:g7:moon-phases:phases}{phase}, a \omterm{def:g7:shadows-eclipses:lunar}{lunar eclipse}
only at another. Name both, justify with the alignment
geometry, and explain why eclipses are nevertheless rare at
those \omterm{def:g7:moon-phases:phases}{phases}.
\end{exercise}

\begin{exercise}[$\star\star$]\label{exo:g7:moon-phases:11}
Astronauts on the near side of the \omterm{def:g2:moon-in-the-sky:moon}{Moon} look up at the Earth.
Describe what ``\omterm{def:g7:moon-phases:phases}{Earth-phases}'' they see through a month, and
what \omterm{def:g7:moon-phases:phases}{phase} the Earth shows them when we see a new moon. (Run
the geometry both ways.)
\end{exercise}

\begin{exercise}[$\star\star\star$]\label{exo:g7:moon-phases:12}
The old \omterm{ex:g2:moon-in-the-sky:shapes}{crescent}'s secret: on clear evenings, the \omterm{ex:g2:moon-in-the-sky:shapes}{crescent}
\omterm{def:g2:moon-in-the-sky:moon}{Moon}'s \emph{dark} part glows faintly --- ``the old moon in the
new moon's arms''. Your grade-two self met this as a riddle;
solve it now completely: trace the \omterm{def:g1:light-and-shadows:source}{light}'s path, explain why
earthshine is brightest precisely at \omterm{ex:g2:moon-in-the-sky:shapes}{crescent} \omterm{def:g7:moon-phases:phases}{phases} (what
\omterm{def:g7:moon-phases:phases}{phase} does the \emph{Earth} show the \omterm{def:g2:moon-in-the-sky:moon}{Moon} then?), and why it
fades toward \omterm{ex:g2:moon-in-the-sky:shapes}{full moon}.
\end{exercise}

\section{Problem: A Month with the Moon}

\begin{problem}[{Weekend problem --- the observatory's open
month; thirty nights of Moon, one wall-chart; the docent's
quiz}]\label{pb:g7:moon-phases:1}
The town observatory runs a public ``Month with the \omterm{def:g2:moon-in-the-sky:moon}{Moon}''. As
junior docent, you keep the wall-chart and answer the visitors.

\medskip\noindent\textbf{Part I --- Setting up the chart.}
\begin{enumerate}
  \item Night 0 is new moon. Fill the chart's milestones: on
        about which nights fall first quarter, \omterm{ex:g2:moon-in-the-sky:shapes}{full moon}, and
        last quarter?
  \item Beside each milestone, the chart wants rise-times.
        Match: rises at \omterm{def:g1:day-and-night:daynight}{sunset}; \ rises near midday; \ rises
        near midnight; \ rises and sets with the Sun.
  \item A visitor asks why the chart shows the \omterm{ex:g2:moon-in-the-sky:shapes}{crescent} low in
        the \emph{western} evening sky but the \omterm{ex:g2:moon-in-the-sky:shapes}{full moon}
        climbing from the \emph{east} at dusk. Explain with
        \omterm{def:g6:sun-earth-moon:axisorbit}{orbit} positions.
  \item Another visitor swears they saw the \omterm{def:g2:moon-in-the-sky:moon}{Moon} at ten in the
        morning. Which chart \omterm{def:g7:moon-phases:phases}{phases} make a morning \omterm{def:g2:moon-in-the-sky:moon}{Moon}
        possible?
\end{enumerate}

\medskip\noindent\textbf{Part II --- The docent's
demonstrations.}
\begin{enumerate}[resume]
  \item You perform the lamp-and-ball walk for a school group.
        A pupil \omterm{def:g2:ice-water-steam:meltfreeze}{freezes} you mid-turn: the ball shows
        three-quarters lit. Name the \omterm{def:g7:moon-phases:phases}{phase}, and locate lamp,
        Earth-you and ball for the group.
  \item ``So the Earth's \omterm{def:g1:light-and-shadows:shadow}{shadow} makes the \omterm{def:g7:moon-phases:phases}{phases}?'' asks the
        pupil, pointing at your dark-sided ball. Give the
        demonstration's own refutation: whose night side are
        they seeing, and where would the Earth's \omterm{def:g1:light-and-shadows:shadow}{shadow}
        actually be?
  \item The group asks why they always see ``the rabbit'' (or
        ``the face'') in the \omterm{def:g2:moon-in-the-sky:moon}{Moon}, never its back. Answer with
        the same-face fact.
  \item A skeptical parent: ``If the \omterm{def:g2:moon-in-the-sky:moon}{Moon} spins, as you claim,
        shouldn't we see it turn?'' Resolve the apparent
        contradiction: how can the \omterm{def:g2:moon-in-the-sky:moon}{Moon} spin exactly once per
        \omterm{def:g6:sun-earth-moon:axisorbit}{orbit} and thereby \emph{hide} its spin from us?
\end{enumerate}

\medskip\noindent\textbf{Part III --- Quiz night.}
\begin{enumerate}[resume]
  \item Quiz card A: ``Tonight's \omterm{def:g2:moon-in-the-sky:moon}{Moon} rose at midnight and
        shows its \emph{left} half lit (as seen going toward
        morning). \omterm{def:g7:moon-phases:phases}{Waxing} or \omterm{def:g7:moon-phases:phases}{waning} --- and which quarter?''
  \item Quiz card B: ``An artist's poster shows a \omterm{ex:g2:moon-in-the-sky:shapes}{crescent} \omterm{def:g2:moon-in-the-sky:moon}{Moon}
        at midnight, high overhead, horns pointing straight
        down at the \omterm{def:g1:day-and-night:daynight}{sunset} colors on the horizon.'' List every
        astronomical error on the poster.
  \item Quiz card C: ``During tonight's \omterm{ex:g2:moon-in-the-sky:shapes}{full moon}, could
        visitors also witness a \omterm{def:g7:shadows-eclipses:solar}{solar eclipse} somewhere on
        Earth?'' Settle it with \omterm{def:g7:moon-phases:phases}{phase} geometry.
  \item Closing speech: in three sentences, give the month's
        one underlying fact, the misconception the month
        buried, and the schedule trick visitors should carry
        home (what any \omterm{def:g7:moon-phases:phases}{phase} tells about its rising hour).
\end{enumerate}
\end{problem}
